% answers_short.tex —— 2011 年·江西卷·理科 · 简短答案速查
% 由 tools/make_short_answers.py 自 answers.tex 的 \gksolhead/\gkanswer 提取，
% 并与 2011/2011_数学真题_江西理_答案要点.md 逐题比对后生成，勿手改（改源 MD 后重跑）。

\begin{gknotice}
\textbf{使用说明}：本节只给最终答案便于快速核对，完整推导见解析版。答案为编写者独立推导，\textbf{非官方标准答案}。
\end{gknotice}

\gkpart{选择题}
\begin{center}\begin{tabular}{ccccccccccc}
\toprule
题号 & 1 & 2 & 3 & 4 & 5 & 6 & 7 & 8 & 9 & 10 \\
\midrule
答案 & D & B & A & C & A & C & D & C & B & A \\
\bottomrule
\end{tabular}\end{center}

\gkpart{填空题}
\begin{itemize}[leftmargin=2.6em]
  \item[\textbf{11.}] $\dfrac\pi3$。
  \item[\textbf{12.}] $\dfrac{13}{16}$。
  \item[\textbf{13.}] $10$。
  \item[\textbf{14.}] $\dfrac{x^2}5+\dfrac{y^2}4=1$。
  \item[\textbf{15.}] A：$(x-2)^2+(y-1)^2=5$；B：$5$。
\end{itemize}

\gkpart{解答题}
\begin{itemize}[leftmargin=2.6em]
  \item[\textbf{16.}] (1) $P(X=k)=\dfrac{\mathrm C_4^k\mathrm C_4^{4-k}}{\mathrm C_8^4}$，$k=0,1,2,3,4$，即 $P(X=0)=\dfrac1{70}$，$P(X=1)=\dfrac{16}{70}$，$P(X=2)=\dfrac{36}{70}$，$P(X=3)=\dfrac{16}{70}$，$P(X=4)=\dfrac1{70}$；(2) $2280$ 元
  \item[\textbf{17.}] (1) $\dfrac34$；(2) $\sqrt7+1$
  \item[\textbf{18.}] (1) $a_n=(2+\sqrt2)^{n-1}$ 或 $a_n=(2-\sqrt2)^{n-1}$；(2) $0<a\leqslant\dfrac13$
  \item[\textbf{19.}] (1) $a>-\dfrac19$；(2) $\dfrac{10}3$
  \item[\textbf{20.}] (1) $\dfrac{\sqrt{30}}5$；(2) $-\dfrac14$
  \item[\textbf{21.}] (1) 作法见解析；(2) $\dfrac{5\sqrt5}3$
\end{itemize}
